\documentclass[12pt,a4paper]{article}

\usepackage{fontawesome5}
\providecommand{\faPencilAlt}{\faIcon{pencil-alt}}

\usepackage{styles/couleurs}
\usepackage{styles/mise_en_page}
\usepackage{styles/boites}
\usepackage{styles/maths}
\usepackage{styles/figures}

\usepackage{enumitem}
\usepackage{array}
\usepackage{tabularx}
\usepackage{amsmath}
\usepackage{amssymb}
\usepackage{graphicx}
\usepackage{tikz}
\usepackage{xparse}

\usetikzlibrary{calc,angles,quotes,arrows.meta}

\setlist[itemize]{leftmargin=1.4em,itemsep=0.15em,topsep=0.2em}
\setlist[enumerate]{leftmargin=*,itemsep=0.22em,topsep=0.25em}

\newcommand{\lignepointillee}{\noindent\dotfill\par\vspace{0.45em}}

\newcommand{\lignesreponse}[1]{%
\begin{tcolorbox}[
enhanced,colback=white,colframe=NeutreC!40,boxrule=0.5pt,
arc=1.5mm,left=3mm,right=3mm,top=2mm,bottom=1mm,
before skip=2pt,after skip=2pt]
\foreach \i in {1,...,#1}{\lignepointillee}
\end{tcolorbox}}

\newcommand{\titreseq}[1]{%
\begin{center}
{\sffamily\large\bfseries\color{SavoirC}#1}
\end{center}}

\begin{document}

\titreseq{Activité 7 -- Mesurer sans pouvoir atteindre}

\section*{Problème 1 -- Le câble de haubanage}

Un câble de \(18\) m relie le sommet d'un mât au sol.
Il forme un angle de \(58^\circ\) avec le sol.

On cherche la distance horizontale entre le pied du mât et le point d'ancrage.

\begin{center}
\begin{tikzpicture}[scale=0.85]
  \coordinate (A) at (0,0);
  \coordinate (B) at (8,0);
  \coordinate (C) at (8,6.4);
  \draw[line width=1pt] (A)--(B)--(C)--cycle;
  \draw (7.6,0)--(7.6,0.4)--(8,0.4);
  \draw pic[draw=ActionC,fill=ActionC!15,angle radius=9mm,"$58^\circ$"] {angle=B--A--C};
  \node[below left] at (A) {ancrage};
  \node[below right] at (B) {pied};
  \node[above right] at (C) {sommet};
\end{tikzpicture}
\end{center}

\begin{enumerate}
  \item Complète la figure.
  \item Quelle longueur est l'hypoténuse ?
  \item Calcule la distance horizontale, arrondie au dixième de mètre.
\end{enumerate}

\lignesreponse{15}

\newpage

\section*{Problème 2 -- Une échelle contre un mur}

Le pied d'une échelle est placé à \(2{,}4\) m du mur.
L'échelle forme un angle de \(68^\circ\) avec le sol.

Quelle longueur minimale doit avoir l'échelle pour atteindre le point visé sur le mur ?

\begin{center}
\begin{tikzpicture}[scale=0.95]
  \coordinate (A) at (0,0);
  \coordinate (B) at (5,0);
  \coordinate (C) at (5,5.9);
  \draw[line width=1pt] (A)--(B)--(C);
  \draw[line width=1pt] (A)--(C);
  \draw (4.65,0)--(4.65,0.35)--(5,0.35);
  \draw pic[draw=ActionC,fill=ActionC!15,angle radius=9mm,"$68^\circ$"] {angle=B--A--C};
  \node[below] at ($(A)!0.5!(B)$) {$2{,}4$ m};
  \node[above left] at ($(A)!0.5!(C)$) {$?$};
\end{tikzpicture}
\end{center}

\lignesreponse{20}

\newpage

\section*{Problème 3 -- La hauteur d'un arbre}

Une observatrice regarde le sommet d'un arbre.
Ses yeux sont à \(1{,}65\) m du sol.
Elle se trouve à \(30\) m horizontalement du tronc.
La ligne de visée forme un angle de \(35^\circ\) avec l'horizontale.

On veut déterminer la hauteur de l'arbre.

\begin{center}
\begin{tikzpicture}[scale=0.27]
  \coordinate (S) at (0,1.65);
  \coordinate (H) at (30,1.65);
  \coordinate (A) at (30,{1.65+30*tan(35)});
  \coordinate (P) at (30,0);
  \draw[line width=1pt] (S)--(H)--(A)--cycle;
  \draw[dashed] (H)--(P);
  \draw[line width=1.2pt] (-2,0)--(34,0);
  \draw (29,1.65)--(29,2.65)--(30,2.65);
  \draw pic[draw=ActionC,fill=ActionC!15,angle radius=18mm,"$35^\circ$"] {angle=H--S--A};
  \node[left] at (S) {$S$};
  \node[right] at (H) {$H$};
  \node[right] at (A) {$A$};
  \node[right] at (P) {$P$};
  \node[below] at ($(S|-P)!0.5!(P)$) {$30$ m};
  \node[left] at ($(S|-P)!0.5!(S)$) {$1{,}65$ m};
\end{tikzpicture}
\end{center}

Le cosinus ne donne pas directement la hauteur \(AH\).
On procède en deux étapes.

\begin{enumerate}
  \item Dans le triangle \(SAH\) rectangle en \(H\), calcule la longueur \(SA\).
  \item À l'aide du théorème de Pythagore, calcule \(AH\).
  \item Déduis la hauteur totale \(AP\) de l'arbre.
  \item Donne le résultat au centimètre près.
\end{enumerate}

\lignesreponse{15}

\end{document}
